
\question
The type of nerve cell responsible for integrating different stimuli is the 

\begin{choices}
\choice sensory neuron 
\choice axon 
\correctchoice interneuron 
\choice motor neuron 
\choice dendrite 
\end{choices}

\question
Transpiration in plants requires all of the following \emph{except}

\begin{choices}
\choice evaporation of water molecules. 
\correctchoice active transport through xylem cells. 
\choice cohesion between water molecules.  
\choice adhesion of water molecules to cellulose. 
\choice water vapor on the spongy mesophyll. 
\end{choices}

\question
Ignoring all other factors, what kind of day would result in the fastest delivery of water and minerals to the leaves of a tree?

\begin{choices}
\choice A cool, dry day. 
\choice A very hot, dry, windy day. 
\correctchoice A warm, dry day. 
\choice A warm, humid day.  
\choice A cool, humid day. 
\end{choices}

\question
Which of the following are primarily responsible for movement in the eukaryotic flagellum?

\begin{choices}
\choice Basal apparatus and hook. 
\correctchoice Dynein and microtubules. 
\choice Actin and myosin filaments.  
\choice Hydrogen ions and basal apparatus. 
\choice Sensory and motor neurons. 
\end{choices}

\question
Hydrostatic skeletons are 

\begin{choices}
\choice external skeletons that provide support and protection against dehydration. 
\correctchoice fluid-filled compartments under pressure. 
\choice skeletons found in deep sea organisms that depend on high pressure of the surrounding water to provide support. 
\choice internal skeletons that are buried in soft tissue.  
\choice skeletons made of microtubules in the cytoplasm of protists. 
\end{choices}

\newpage

\question
Galvanotaxis occurs when 

\begin{choices}
\choice protists respond to touch. 
\choice animals respond to temperature. 
\correctchoice protists respond to an electrical field. 
\choice plants respond to gravity.  
\choice plants respond to hormones. 
\choice animals respond to metal ions. 
\end{choices}

\question
Pores on the leaf surface that function in gas and water exchange are called

\begin{choices}
\correctchoice stomata. 
\choice waxy cuticles. 
\choice spongy mesophyll. 
\choice phloem. 
\choice xylem. 
\end{choices}

\question
In terms of structure, cilia  are most like 
\begin{choices}
\choice microtubules.  
\choice myosin filaments. 
\choice actin filaments. 
\choice bacterial flagella. 
\correctchoice eukaryotic flagella. 
\end{choices}

\question
The ion used most often as a second messenger is

\begin{choices}
\choice \ce{K+} 
\choice \ce{H+} 
\choice \ce{Na+} 
\choice \ce{Cl-} 
\correctchoice \ce{Ca+} 
\end{choices}

\question
The fluid that moves around in the circulatory system of a typical arthropod is the

\begin{choices}
\choice cytoplasm. 
\choice blood plasma. 
\correctchoice hemolymph. 
\choice intracellular fluid. 
\choice gastrovascular fluid. 
\end{choices}

\question
Most of the water taken up by a plant

\begin{choices}
\choice transports sugars in the phloem. 
\choice keep guard cells turgid. 
\correctchoice is lost during transpiration. 
\choice is converted to sugars during photosynthesis. 
\choice enters the central vacuoles of cells.  
\end{choices}

\question
Organisms with a circulating fluid that is separate from the fluid that directly surrounding the body's cells are likely to have

\begin{choices}
\choice a hemolymph system. 
\choice a vascular system. 
\correctchoice a closed circulatory system. 
\choice an open circulatory system. 
\choice a gastrovascular cavity. 
\end{choices}

\question
The main stimulus used by birds to find their way north for the summer time is 
\begin{choices}
\correctchoice electomagnetic.
\choice thigmic. 
\choice chemical. 
\choice thermal. 
\choice mechanical. 
\end{choices}

\question
What drives the flow of water through the xylem?

\begin{choices}
\choice The number of companion cells in the phloem.  
\choice Abscisic acid opening guard cells around the stomata. 
\choice Transpirational pull through sieve-tube elements. 
\correctchoice The evaporation of water from the leaves. 
\choice A higher concentration of sugars in the leaves than in the roots. 
\end{choices}

\question
During protist conjugation, which of the following does \emph{not} occur during some part of the process?

\begin{choices}
\choice Some micronuclei distintegrate. 
\correctchoice The macronucleus fuses with the micronucleus. 
\choice The micronucleus undergoes mitosis. 
\choice The macronucleus distintegrates. 
\choice Cells exchange micronuclei.  
\choice The micronucleus undergoes meiosis. 
\end{choices}


\question
From an energy perspective, animals are mobile because they consume

\begin{choices}
\choice low quality energy and low energy efficiency.
\choice high quality energy and carbon intake is irregular. 
\choice low quality energy and carbon intake is constant. 
\correctchoice low quality energy and carbon intake is irregular.
\choice high quality energy and carbon intake is constant. 
\choice high quality energy and high energy efficiency. 
\end{choices}

\newpage

\question
The highest quality form of energy (best able to do work) for organisms is

\begin{choices}
\choice thermal.
\correctchoice radiant. 
\choice chemical. 
\choice potential. 
\choice kinetic. 
\end{choices}

\question
When ATP attaches to the myosin head

\begin{choices}
\choice the actin head attaches to the myosin filament. 
\correctchoice the myosin head releases from the actin filament. 
\choice the myosin head initiates the power stroke. 
\choice the actin head attaches to the myosin filament. 
\choice the myosin head attaches to the actin filament.  
\end{choices}


\question
The protein that is activated in response to auxin that allows phototropism to occur is

\begin{choices}
\choice transmembrane. 
\choice wallulase. 
\choice basal apparatus. 
\choice abscisic acid. 
\correctchoice expansin. 
\choice proton channel. 
\end{choices}

\question
The vascular tissue system for delivery of water, minerals, and nutrients throughout the plant includes all of the following except for

\begin{choices}
\choice phloem. 
\choice xylem. 
\correctchoice cohesion. 
\choice companion cells. 
\choice sieve plates. 
\end{choices}

\question
Which of the following about phloem is false?

\begin{choices}
\choice Sugars move by passive diffusion. 
\choice Companion cells provide proteins to sieve-tube elements. 
\choice The cell has cytoplasm. 
\choice Sieve-tube elements do not have a nucleus. 
\correctchoice The plasma membrane is modified to form a sieve plate. 
\end{choices}

\newpage

\question
Differences in energy quality and availability affects the amount of surface area and volume in organisms. Compared to animals, plants have 

\begin{choices}
\choice low surface area and high volumes. 
\correctchoice high surface area and low volumes. 
\choice high surface area and high volumes. 
\choice low surface areas and low volumes.  
\end{choices}


\question
The function of valves in veins in a closed circulatory system is to 
\begin{choices}
\choice enhance diffusion of nutrients to the tissues. 
\correctchoice ensure one way flow of circulatory fluids. 
\choice enable the vein to withstand high pressure. 
\choice provide support for the extra muscles around the vein. 
\choice decrease the amount of pressure in the veins.  
\end{choices}
